% ============================================================
% Section 03 — Invariance and Continuity
% Book I: Ontology of Viability
% Chapter 05: Identity and Continuity
%
% Purpose:
%   To examine the ontological relation between invariance
%   and continuity without reducing identity continuity
%   to the preservation of invariant features.
%
% Conceptual Policy:
%   1. Invariance is not identity.
%   2. Invariance is not persistence or continuity.
%   3. The preservation of an invariant does not by itself
%      establish the continuity of identity.
%   4. The absence of an identified invariant does not
%      by itself establish discontinuity of identity.
%   5. Every invariant must be specified relative to a
%      transformation, domain, and respect.
%   6. Formal invariance must not automatically be treated
%      as ontological invariance.
% ============================================================

\section{Invariance and Continuity}
\label{sec:ch05-invariance-and-continuity}

\subsection{The Ontological Problem of Invariance and Continuity}

The relation between invariance and continuity constitutes
a central problem in the investigation of identity
through change.

An existent may undergo transformations while preserving
certain features, relations, structures, or determinations.
Such preservation may be described in terms of invariance.
The existence of invariants can therefore provide relevant
information about the conditions under which an existent
continues through transformation.

However, the preservation of an invariant and the
continuity of identity are not equivalent ontological
conditions.

An invariant may remain preserved across a specified
transformation without establishing that the numerically
same identity has continued. Conversely, an existent
may undergo transformations in which no previously
identified invariant remains preserved, without this
fact alone establishing that its identity has ceased.

The problem is therefore not simply whether something
remains invariant, but whether the invariant in question
has a justified ontological relation to the identity
whose continuity is being examined.

This requires distinguishing the formal preservation
of a feature from the ontological preservation of
identity. It also requires identifying the transformations
under consideration and the domain and respect in which
invariance is asserted.

The present section examines these relations without
assuming that identity is constituted by invariance
or that continuity can be established through the
preservation of invariant features alone.

\subsection{Invariance as a Relation to Transformation}

Invariance is meaningful only in relation to a specified
class of transformations.

A feature, relation, structure, or determination is
invariant when it remains preserved under the relevant
transformations according to an explicitly defined
criterion.

Consequently, an assertion of invariance requires
at least three elements:

\begin{itemize}
    \item a feature or property whose preservation
    is being examined;
    \item a transformation or class of transformations
    relative to which preservation is assessed;
    \item a domain and respect in which the invariance
    is asserted.
\end{itemize}

Without these specifications, the claim that something
is invariant remains incomplete.

A feature may remain invariant under one class of
transformations while changing under another. A relation
may be preserved in one domain while failing to remain
preserved in a different domain. A structure may remain
formally equivalent under a transformation even though
its material realization changes.

Invariance is therefore not an absolute description
of an existent independently of the transformations
and respects under consideration.

This relativity to transformation is particularly
important when invariance is invoked in discussions
of identity continuity. The preservation of a feature
under a limited class of transformations cannot be
extended without justification to all possible
transformations of the existent.

\subsection{The Distinction Between Invariance and Continuity}

Invariance and continuity address different questions.

Invariance concerns the preservation of a specified
feature under a specified transformation.

Continuity concerns whether a relevant existent,
identity, process, or manifestation continues across
a specified succession or domain of change.

The first is therefore a relation between a feature
and a transformation. The second concerns the continuation
of a subject through a relevant domain.

An invariant may be preserved across a transformation
without establishing the continuity of the existent
to which it is attributed.

Similarly, continuity may obtain in a relevant sense
even when some features undergo alteration and no
single previously identified invariant remains preserved.

The two notions may be related, but their relation
must be established rather than assumed.

The preservation of an invariant can contribute to
an account of continuity when the invariant has been
shown to be ontologically relevant to the subject
under examination. Yet the mere identification of
a preserved feature does not establish this relevance.

Likewise, the transformation of a feature does not
automatically establish discontinuity. It may indicate
a change in a particular determination or respect
without determining whether the identity of the
existent has continued.

\subsection{Invariance and the Persistence of Identity}

The persistence of identity must not be defined as
the preservation of an invariant.

Such a definition would make the ontological question
of identity dependent upon a criterion that has not
yet been independently justified.

If identity were defined through the preservation
of a selected invariant, the subsequent claim that
identity persists because the invariant remains
preserved would merely repeat the initial definition.
It would not establish an independent ontological
account of identity continuity.

Furthermore, the selection of an invariant requires
an explanation of why that particular feature should
be regarded as relevant to identity.

A material feature, structural relation, functional
organization, or relational configuration may remain
preserved across transformations. However, the preservation
of any such feature does not automatically establish
that the identity associated with it has persisted.

The same invariant may characterize distinct existents.
Consequently, the preservation of a feature cannot
by itself establish numerical identity between the
subjects that exhibit it.

Conversely, the transformation or disappearance of
a particular invariant does not independently establish
that the identity of its bearer has ceased.

The relation between invariance and identity persistence
must therefore be investigated without making either
notion a substitute for the other.

\subsection{The Problem of Selecting Identity-Relevant Invariants}

Not every invariant has the same ontological significance.

A feature may remain preserved because it belongs
to a relatively stable configuration, because the
considered transformation does not affect it, or
because the criterion of invariance is defined at
a level of abstraction that excludes certain changes.

None of these circumstances independently establishes
that the feature is constitutive of identity.

The identification of an identity-relevant invariant
therefore requires more than demonstrating preservation.

It requires an account of the relation between the
invariant and the existent whose identity is under
examination.

At least three questions must be distinguished:

\begin{enumerate}
    \item Is the feature genuinely invariant under
    the transformations being considered?
    \item Is the feature ontologically relevant to
    the identity of the existent?
    \item Does the preservation of this feature provide
    sufficient grounds for establishing continuity
    of identity?
\end{enumerate}

An affirmative answer to the first question does
not entail affirmative answers to the second or third.

A feature may be invariant without being identity-relevant.
It may be identity-relevant without being sufficient
to establish identity continuity independently.
It may also contribute to an account of identity
preservation only in conjunction with other conditions.

The ontological role of an invariant must therefore
be established for the particular existent, domain,
and class of transformations under examination.

\subsection{The Non-Equivalence of Formal and Ontological Invariance}

Formal invariance and ontological invariance must
remain distinct.

A formal system may identify quantities, relations,
or structures that remain unchanged under a specified
transformation. Such invariance is established relative
to the definitions, rules, and domain of the formal
system.

Ontological invariance concerns the preservation of
a feature or determination in the reality of an existent
under the relevant transformations.

Formal invariance may provide a useful means of
representing or investigating ontological relations.
Nevertheless, the preservation of a mathematical
or formal structure does not automatically establish
that the corresponding ontological feature has been
preserved.

A formal mapping may preserve a relation while the
entities represented by that relation undergo changes
whose ontological significance is not captured by
the mapping.

Likewise, formal equivalence between two configurations
does not independently establish that they belong
to one and the same numerical identity.

The transition from formal invariance to ontological
invariance therefore requires an argument establishing
the adequacy of the formal representation for the
ontological subject under examination.

Without such an argument, formal preservation remains
a result internal to the formal domain.

\subsection{Transformation Classes and the Scope of Continuity}

The relation between invariance and continuity depends
upon the class of transformations being considered.

An existent may preserve a feature under minor structural
alterations while failing to preserve it under more
extensive transformations. A relation may remain
invariant under transformations of material realization
but not under transformations of relational organization.

The continuity of identity cannot therefore be inferred
from invariance under a restricted class of transformations
unless the relevance and adequacy of that class have
been established.

The distinction also prevents the unrestricted use
of the proposition that identity remains invariant
under all temporal transformations.

Temporal transformations do not constitute a
self-justifying class of transformations under which
identity must necessarily remain invariant.

The fact that a transformation occurs through time
does not establish that it is compatible with the
persistence of identity. Nor does it establish that
it is incompatible with such persistence.

The ontological status of each transformation must
be examined independently.

A claim of continuity through transformation must
therefore specify which transformations are involved,
which features remain preserved, and why those features
are relevant to the identity under examination.

\subsection{Continuity Without an Identified Invariant}

The absence of an identified invariant does not
necessarily establish the discontinuity of identity.

There may be several reasons why an invariant has
not been identified.

The relevant feature may not yet be known. The
transformation class may not have been adequately
specified. The chosen level of analysis may fail
to capture the relevant ontological relation.
Alternatively, identity continuity may not depend
upon the preservation of a single invariant feature.

Consequently, the failure to identify an invariant
must be distinguished from the demonstrated absence
of every possible identity-relevant invariant.

Even the demonstrated transformation of all previously
identified features does not automatically establish
that the numerical identity of the existent has ceased.

This does not imply that identity must persist whenever
no invariant is identified. It means only that the
absence of an identified invariant is insufficient
as an independent proof of discontinuity.

The question remains open until the relevant ontological
conditions have been examined.

\subsection{Invariance Without Continuity of Identity}

The reverse possibility must also be considered.

An invariant may remain preserved across a transformation
even when the numerical identity of the subject
has not been established as continuous.

Distinct existents may share the same invariant feature.
A structure may be reproduced in another existent.
A relation may be maintained across a transition
between different subjects.

In these circumstances, the preservation of the
invariant establishes a form of continuity or
equivalence at the level of the relevant feature,
but it does not necessarily establish continuity
of numerical identity.

This possibility demonstrates why invariance cannot
serve as a sufficient criterion of identity persistence.

The ontological question concerns not merely whether
a feature remains the same, but whether the existent
whose identity is being examined remains numerically
the same existent.

The distinction between feature preservation and
identity preservation must therefore be maintained
even when the invariant is highly stable or structurally
significant.

\subsection{Invariance, Change, and the Limits of Preservation}

Invariance must not be understood as the absence
of all change.

An existent may preserve certain features while
undergoing transformations in others. The preservation
of one respect does not imply the immutability of
the whole existent.

Likewise, the alteration of a feature does not
necessarily establish that every relevant form of
continuity has been interrupted.

The ontological significance of change depends upon
the subject and respect under examination.

A transformation may preserve an invariant at one
level while altering a different determination at
another. It may preserve a structural relation while
changing its material realization. It may preserve
a relational organization while altering the functions
performed within it.

These possibilities demonstrate that invariance and
change are not mutually exclusive descriptions
of an existent considered without qualification.

The preservation of selected features must not be
generalized into a claim of total preservation.
Likewise, the alteration of selected features must
not be generalized into a claim of total discontinuity.

An adequate account of identity continuity must
therefore specify the respects in which preservation
and transformation occur.

\subsection{The Relation Between Invariance and Manifestation Continuity}

The continuity of manifestations may involve the
preservation of certain invariant features.

A sequence of manifestations may preserve a structural
relation, a material characteristic, or another
specified determination across successive conditions.

Such preservation may help describe the continuity
of the manifestation sequence. Nevertheless, it
does not independently establish the persistence
of the identity associated with those manifestations.

The same invariant feature may occur in manifestations
belonging to distinct existents. Conversely, manifestations
associated with one existent may undergo transformations
that alter previously preserved features.

The relation between manifestation continuity and
identity continuity must therefore be established
independently of the mere preservation of manifest
features.

The continuity of a manifestation sequence may be
relevant to the investigation of identity, but it
cannot be treated as a direct substitute for it.

\subsection{The Subject and Respect of Invariance}

Every claim concerning invariance and continuity
must identify its subject and respect.

The subject may be an existent, an identity, a
determination, a structure, a relation, a material
configuration, or a manifestation.

The respect must specify the feature whose preservation
is being examined and the transformations under
which that preservation is asserted.

An invariant identified at the level of structure
cannot automatically be transferred to the level
of identity. Nor can the preservation of a feature
in one domain be generalized to every domain in
which the existent undergoes transformation.

The same discipline applies to continuity. A claim
that continuity has been established must identify
what continues and in which respect.

Without these specifications, the relation between
invariance and continuity remains ambiguous and
cannot support a determinate ontological conclusion.

\subsection{Methodological Requirements}

The investigation of invariance and continuity must
observe the following methodological requirements:

\begin{enumerate}
    \item Invariance must be specified relative to
    an identifiable transformation or class of
    transformations.
    \item The domain and respect of every invariance
    claim must be explicitly identified.
    \item Formal invariance must not automatically
    be treated as ontological invariance.
    \item The preservation of an invariant must not
    be treated as sufficient proof of identity continuity.
    \item The absence of an identified invariant must
    not be treated as sufficient proof of identity
    discontinuity.
    \item The preservation of a feature must be
    distinguished from the preservation of numerical
    identity.
    \item The relation between invariance and identity
    must be established independently rather than
    assumed through definition.
\end{enumerate}

These requirements prevent invariance from becoming
an unrestricted criterion through which every question
of identity continuity is prematurely resolved.

\subsection{Scope of the Present Section}

This section has examined the relation between
invariance and continuity while preserving their
ontological distinction.

Invariance concerns the preservation of a specified
feature under a specified class of transformations.
Continuity concerns the continuation of a subject
across a relevant domain or succession. Neither
notion can be reduced to the other.

The preservation of an invariant does not independently
establish continuity of identity, and the absence
of an identified invariant does not independently
establish discontinuity of identity.

The ontological relevance of invariants must be
demonstrated in relation to the identity under
examination, the transformations involved, and
the domain and respect of the claim.

The next section examines identity and non-identity
through difference. It investigates whether qualitative,
structural, material, or relational differences
between successive conditions can establish numerical
non-identity, and why difference must not be treated
as an immediate proof of the termination or replacement
of identity.

% ============================================================
% End of Section 03
% ============================================================